% Variante de non-regression — en-tete article, etablissement inp seul.
% Complement de header-article.tex (institution={uftmp,n7}) : le cas simple
% d'une valeur unique autre que le defaut n7.
\documentclass[11pt]{ocots-td}
\usepackage[lang=fr, theme=ocots, solutions=none, math=analysis,
            institution=inp]{ocots}
\lstset{language=[LaTeX]TeX}

\title{TD de démonstration — En-tête INP}
\shorttitle{TD — En-tête}
\numero{TD~o 1}
\date{2025--2026}
\discipline{Calcul différentiel}
\promotion{Toulouse INP — 1A}

\newcommand{\ocotsvariantnote}[1]{\begin{quote}\itshape Variante : #1\end{quote}}

\begin{document}

\maketitle

\begin{lstlisting}
\usepackage[..., institution=inp]{ocots}
\end{lstlisting}
\ocotsvariantnote{\texttt{institution=inp} (le défaut est
\texttt{institution=n7}) : une seule valeur — le logo Toulouse INP occupe
seul le côté gauche de l'en-tête.}

\begin{instructions}
    Le logo de la première page est composé en mode en-tête.
\end{instructions}

\begin{exercise}[points=2]
    Un exercice, pour vérifier que le contenu suit l'en-tête.
\end{exercise}

\end{document}
